\documentclass[12pt]{article}
\usepackage[utf8]{inputenc}
\usepackage[T1]{fontenc}
\usepackage{hyperref}
\title{Sere Language Resources for K-12 Curriculum Development}
\author{Compiled by Claude Sonnet 5.5 (claude-sonnet-5-5) for Andr\'as Kornai}
\date{}
\begin{document}
\maketitle

\section{Introduction}
Sere (also Shere, Siri, Basiri) is a Ubangian language of northeastern DRC (Bas-Uele) and the Central African Republic. Glottolog rates it nearly extinct, and Wikitongues lists it as unwritten. Speaker counts in the sources conflict: about 2,500 (an old figure), 30 in CAR in 1988, and 14,000 for the ethnic group per Joshua Project. Several unrelated languages share the name ``Ngbaka'', so care is needed when searching.

\section{Primary Educational Resources}
None found. I found no evidence of Sere schooling, primers or an orthography.

\section{Linguistic and Reference Materials}
\begin{itemize}
\item Wikipedia: \url{https://en.wikipedia.org/wiki/Sere_language}
\item Glottolog entry: \url{https://glottolog.org/resource/languoid/id/sere1264}
\item Numerals page (Max Planck, Leipzig): \url{https://lingweb.eva.mpg.de/channumerals/Sere.htm}
\item Ethnologue subgroup page: \url{https://www.ethnologue.com/subgroup/1029/}
\end{itemize}

\section{Religious and Community Resources}
Joshua Project people-group profile: \url{https://joshuaproject.net/people_groups/14865}

\section{Digital and Audio Resources}
Wikitongues language page (no Sere recordings were confirmed in the search results): \url{https://wikitongues.org/languages/swf/}

\section{Curriculum Development Context}
Practically no material exists in or about Sere beyond basic descriptions. A K-12 curriculum would depend on neighboring languages of wider communication, and Sere itself would need documentation fieldwork first.

\end{document}